% !TeX spellcheck = ru_RU
% !TEX root = vkr.tex

\section{Заключение}
\label{sec:conclusion}

В ходе работы была реорганизована инфраструктура для измерения производительности графовых алгоритмов на основе \INs{}: вместо набора разрозненных скриптов создана единая система с общим интерфейсом запуска.

В результате решения поставленных задач получены следующие результаты:

\begin{enumerate}
    \item Все этапы работы объединены в одну программу с командами получения матриц, их подготовки и запуска измерений, а также с общим журналом работы.
    \item Общая логика~--- выбор матриц, повторение запусков, перебор числа потоков и сохранение результатов~--- вынесена в отдельный модуль, поэтому код конкретного инструмента содержит только его особенности.
    \item Состав измерений полностью задаётся одним файлом настроек, а загрузка и предобработка матриц, включая построение переупорядоченных версий, выполняются автоматически.
    \item Описаны единые правила подключения инструментов, благодаря чему добавление нового инструмента сводится к нескольким предсказуемым шагам.
    \item К новой инфраструктуре подключены все используемые инструменты, что подтверждает её работоспособность: измерения выполняются через единый интерфейс, результаты сохраняются в общем виде.
\end{enumerate}

Реорганизованная инфраструктура сокращает дублирование кода, упрощает добавление новых инструментов и матриц, а также делает проведение и повторение измерений более удобным.
